% Part: second-order-logic
% Chapter: sol-and-sets
% Section: comparing-sets

\documentclass[../../../include/open-logic-section]{subfiles}

\begin{document}

\olfileid{sol}{set}{cmp}

\olsection{比较集合}

\begin{prop}
公式 $\lforall[x][(X(x) \lif Y(x))]$ 定义了子集关系，即 $\Sat{M}{\lforall[x][(X(x) \lif Y(x))]}[s]$ 当且仅当 $s(X)
\subseteq s(Y)$。
\end{prop}

\begin{prop}
公式 $\lforall[x][(X(x) \liff Y(x))]$ 定义了集合上的恒等关系，即 $\Sat{M}{\lforall[x][(X(x) \liff Y(x))]}[s]$
当且仅当 $s(X) = s(Y)$。
\end{prop}

\begin{prop}
公式 $\lexists[x][X(x)]$ 定义了非空性，即 $\Sat{M}{\lexists[x][X(x)]}[s]$ 当且仅当 $s(X) \neq
\emptyset$。
\end{prop}

一个集合~$X$ 不大于集合~$Y$，记作 $\cardle{X}{Y}$，当且仅当存在一个单射函数 $f\colon X \to Y$。由于我们可以表达一个函数是单射的，并且它在~$X$ 中的自变元上的取值都在~$Y$ 中，因此我们也能在论域的子集上定义“不大于”这一关系。

\begin{prop}
公式
\[
\lexists[u][(\lforall[x][(X(x) \lif Y(u(x)))] \land \lforall[x][\lforall[y][(\eq[u(x)][u(y)] \lif
      \eq[x][y])]])]
\]
定义了“不大于”关系。
\end{prop}

两个集合大小相同，或称“等势”，记作 $\cardeq{X}{Y}$，当且仅当存在一个双射函数~$f\colon X \to Y$。

\begin{prop}
公式
\begin{multline*}
  \lexists[u][(\lforall[x][(X(x) \lif Y(u(x)))] \land {}\\
    \lforall[x][\lforall[y][(\eq[u(x)][u(y)] \lif \eq[x][y])]]
  \land {} \\\lforall[y][(Y(y) \lif \lexists[x][(X(x)
      \land \eq[y][u(x)])])])]
\end{multline*}
定义了“与……等势”关系。
\end{prop}

我们将这些公式分别缩写为 $X \subseteq Y$、
$X = Y$、$X \neq \emptyset$、$\cardle{X}{Y}$ 和
$\cardeq{X}{Y}$。（这可能会引起一些混淆，因为当我们非形式地谈论集合 $X$ 和 $Y$ 时也使用相同的记号——但在这里，这些记号是涉及一元关系变元 $X$ 和~$Y$ 的二阶逻辑公式的缩写。）

\begin{prop}
语句 $\lforall[X][\lforall[Y][((\cardle{X}{Y} \land
    \cardle{Y}{X}) \lif \cardeq{X}{Y})]]$ 是有效的。
\end{prop}

\begin{proof}
该语句在一个结构~$\Struct{M}$ 中被满足，如果对任意子集 $X \subseteq \Domain{M}$ 和 $Y \subseteq \Domain{M}$，若 $\cardle{X}{Y}$ 且 $\cardle{Y}{X}$，则 $\cardeq{X}{Y}$。但这对于\emph{任何}集合 $X$ 和 $Y$ 都成立——它正是 Schr\"oder-Bernstein 定理。
\end{proof}


\end{document}